\documentclass[a4paper, 12pt]{article}
\usepackage[utf8]{inputenc}
\usepackage{mathptmx}
\usepackage[T1]{fontenc}
\usepackage{listingsutf8}
\usepackage{graphicx}
\usepackage{longtable}
\usepackage{booktabs}
\usepackage{listings}
\usepackage{xcolor}
\usepackage{enumitem}

\definecolor{codegreen}{rgb}{0,0.6,0}
\definecolor{codegray}{rgb}{0.5,0.5,0.5}
\definecolor{codepurple}{rgb}{0.58,0,0.82}
\definecolor{backcolour}{rgb}{0.95,0.95,0.92}

\lstset{
    backgroundcolor=\color{gray!10},   
    commentstyle=\color{codegreen},
    keywordstyle=\color{magenta},
    numberstyle=\tiny\color{codegray},
    stringstyle=\color{codepurple},
    basicstyle=\ttfamily\footnotesize,
    breakatwhitespace=false,         
    breaklines=true,                 
    captionpos=b,                    
    keepspaces=true,                 
    numbers=left,                    
    numbersep=5pt,                  
    showspaces=false,                
    showstringspaces=false,
    showtabs=false,                  
    tabsize=2,
    inputencoding=utf8,
    extendedchars=false,
    escapeinside={\%*}{*)},
}
\usepackage{hyperref}
\hypersetup{
	pdftitle={When Blocks Build Societies: Multi-Agent Systems in Minecraft},
    pdfauthor={Jan-Philipp Kiel},
	pdfnewwindow=true,
	% colorlinks=false,
	% linktoc=page,     % nur seitennummern verlinkt
	linktoc=all, % sections und subsections verlinkt
	% linkcolor=black,
	% urlcolor=black,
	% bookmarks=true,
	% citecolor=black
}
\usepackage[labelfont=bf]{caption}
\usepackage{float}

% set margins for double-sided printing
\usepackage[left=3.5cm, right=2.5cm, top=2.5cm, bottom=2.5cm]{geometry} 
\usepackage{setspace}
\onehalfspacing
% set header
\usepackage{fancyhdr}
\pagestyle{fancy}
\fancyhead{}
\fancyfoot{}
\fancyhead[L]{When Blocks Build Societies: Multi-Agent Systems in Minecraft}
\fancyfoot[R]{\thepage}
\fancypagestyle{plain}[fancy]{} % use fancy style for TOC
\setlength{\headheight}{15pt}
\usepackage{indentfirst} % indent first paragraph after section heading

\usepackage{tocloft}
\renewcommand{\cftsecleader}{\cftdotfill{\cftdotsep}}

% set APA citation style
\usepackage{apacite}
\usepackage[nottoc]{tocbibind}

% Version Watermark
\usepackage{gitinfo2} % Gitinfo
\newcommand{\lastChange}{\gitCommitterIsoDate}
\usepackage{background}
\backgroundsetup{
	contents=\Large{\texttt{\MakeUppercase{\gitBranch}-\gitDescribe{ @ }\lastChange}},
	color=red,
	placement=center,
	angle=90,
	opacity=1,
	vshift=9cm,
	scale=1
}

\begin{document}

\begin{titlepage}
    \begin{center}
        \vspace*{\fill}
        \vspace*{-2cm}
        \large
        {\bfseries{When Blocks Build Societies}}\\
        Multi-Agent Systems in Minecraft\\

        \vspace{1cm}

        Jan-Philipp Kiel\\
        University of Cologne

        \vspace{1cm}

        Advanced Seminar\\
        Information Systems and Digital Technology\\
        Prof. Dr. Stefan Seidel

        \vspace{1cm}

        \today
        \vspace*{\fill}
    \end{center}
\end{titlepage}

\clearpage
\pagenumbering{Roman}
\tableofcontents
\clearpage
% \listoffigures
% \clearpage
% \listoftables
% \clearpage

\pagenumbering{arabic}


%%%%%%%%%%%%%%%%%%%%%%%%%%%%%%%%%%%%%%%%%%%%%%%%%%%%%%%%%%%%%
%MAIN PART
%%%%%%%%%%%%%%%%%%%%%%%%%%%%%%%%%%%%%%%%%%%%%%%%%%%%%%%%%%%%%

\section{Introduction}

Multi-agent systems represent one of the most promising frontiers in artificial intelligence (AI) research, offering insights into how autonomous entities can collaborate, compete, and coexist within shared environments. These systems enable the study of complex emergent behaviors and collective intelligence that extend beyond the capabilities of single agents \cite{Leeuwen2021SelfOrganizingMS, shu2024effectivegenaimultiagentcollaboration}. As AI continues to advance, understanding the principles that govern effective multi-agent collaboration becomes increasingly vital for applications ranging from robotic teams to virtual assistants.

Minecraft, with its open-world sandbox nature, procedurally generated environments, and flexible crafting systems, has emerged as an ideal platform for studying multi-agent interactions \cite{fan2022minedojobuildingopenendedembodied}. The game's combination of simple rules and complex emergent behaviors creates a rich testbed for AI research that bridges the gap between controlled laboratory experiments and the messy complexity of real-world environments \cite{zhao2024steveseriesstepbystepconstruction}. Its block-based world allows for precise control while still presenting meaningful challenges in navigation, resource gathering, construction, and social interaction.

Current multi-agent research in Minecraft has predominantly focused on specific task completion, such as navigation or building predefined structures. However, a significant gap exists between these narrowly defined task-oriented systems and the more complex socially-oriented agent societies that could potentially emerge \cite{shu2024effectivegenaimultiagentcollaboration, qiu2024collaborativeintelligencepropagatingintentions}. Bridging this gap requires new approaches to agent architecture, communication protocols, and coordination mechanisms that can support sustainable collaboration over extended periods.

Two recent research initiatives have made significant strides in addressing different aspects of this challenge: Project Sid and TeamCraft. Project Sid, introduced in 2024, explores large-scale simulations of 10-1000+ AI agents, focusing on how these agents develop specialized roles, adhere to collective rules, and engage in cultural transmission within a civilization-inspired framework \cite{al2024projectsidmanyagentsimulations}. In contrast, TeamCraft provides a multi-modal multi-agent benchmark featuring 55,000 task variants designed to test various aspects of agent collaboration through structured prompts and demonstrations \cite{long2024teamcraftbenchmarkmultimodalmultiagent}.

This paper addresses the research question: 

\begin{quote}
    \emph{How do Project Sid's self-organizing social systems compare to TeamCraft's task-oriented coordination frameworks in achieving sustainable multi-agent collaboration?}
\end{quote}

This comparison is crucial for understanding the trade-offs between emergent, socially-driven organization and structured, task-driven coordination in multi-agent systems. By analyzing these contrasting approaches, we aim to identify design principles that can inform the development of more robust and adaptable collaborative AI systems.

The paper is structured as follows: First, we review background literature on multi-agent systems, Minecraft as a research platform, and approaches to agent collaboration. Next, we detail our methodology, including case study selection criteria and analytical framework. We then analyze Project Sid and TeamCraft individually before conducting a comparative analysis across organizational structures, communication patterns, learning mechanisms, and sustainability. Finally, we discuss theoretical implications and practical applications before concluding with directions for future research.

\section{Background and Related Work}

\subsection{Multi-Agent Systems in Virtual Environments}

Multi-agent systems represent a paradigm in artificial intelligence research where multiple autonomous agents interact within a shared environment to achieve individual or collective goals. The field has evolved significantly from early distributed AI research in the 1980s to contemporary approaches leveraging deep learning and reinforcement learning \cite{Leeuwen2021SelfOrganizingMS, shu2024effectivegenaimultiagentcollaboration}. Modern multi-agent systems research encompasses diverse architectures ranging from fully cooperative to competitive and mixed-motive scenarios, with applications spanning robotics, logistics, gaming, and social simulation.

Virtual environments offer particularly valuable platforms for multi-agent research due to their controlled conditions, cost-effectiveness, and ability to simulate complex scenarios that would be impractical or impossible in physical settings \cite{shu2024effectivegenaimultiagentcollaboration, Lhaksmana2018RoleBasedMF}. These environments provide a "consequence-free" space to test agent behaviors that might be risky or expensive in real-world settings, while still capturing essential dynamics of multi-agent interaction. Additionally, virtual environments allow researchers to precisely control experimental conditions and collect comprehensive data about agent behaviors and interactions.

Multi-agent architectures can be broadly categorized as centralized or decentralized, with various hybrid approaches emerging to leverage the advantages of both \cite{Leeuwen2021SelfOrganizingMS}. Centralized approaches maintain a global view of the environment and coordinate agent actions through a central controller, offering optimal decision-making at the cost of scalability and robustness. Decentralized approaches distribute decision-making among individual agents, providing resilience and scalability at the expense of potentially suboptimal coordination. Communication models within these architectures range from direct message passing to stigmergic coordination through environmental modifications.

Recent advances in multi-agent collaboration have been driven by developments in large language models (LLMs) and deep reinforcement learning \cite{shu2024effectivegenaimultiagentcollaboration, qiu2024collaborativeintelligencepropagatingintentions}. LLMs have enhanced agents' ability to understand complex instructions and generate contextually appropriate responses, while reinforcement learning techniques have improved adaptation to dynamic environments and coordination among heterogeneous agents \cite{chen2024sagentsselforganizingagentsopenended}. These advances have expanded the scope of multi-agent research to include emergent communication protocols, collective intelligence, and social learning.

\subsection{Minecraft as a Research Platform}

Minecraft offers a uniquely suitable platform for AI research due to its procedurally generated worlds, comprehensive crafting system, and physics simulation that strikes a balance between realism and abstraction \cite{fan2022minedojobuildingopenendedembodied}. The game's block-based structure provides a discretized state space that is more computationally tractable than continuous environments while still offering sufficient complexity to model interesting real-world problems. Furthermore, its open-world nature allows for a wide range of tasks, from simple navigation to complex collaborative construction projects.

Previous AI research in Minecraft has explored diverse challenges, including single-agent task completion, navigation in procedurally generated landscapes, and creative building \cite{zhao2024steveseriesstepbystepconstruction, wang2023jarvis1openworldmultitaskagents}. The MineDojo framework, for instance, provides thousands of diverse open-ended tasks and an internet-scale knowledge base with Minecraft videos, tutorials, and forum discussions to facilitate agent learning \cite{fan2022minedojobuildingopenendedembodied}. Similarly, JARVIS-1 demonstrates how multimodal language models can perceive visual observations and textual instructions to generate sophisticated plans and perform embodied control in Minecraft \cite{wang2023jarvis1openworldmultitaskagents}.

The Minecraft environment presents unique challenges for AI research, including a large state space, partial observability, and the need for long-horizon planning \cite{fan2022minedojobuildingopenendedembodied, wang2023jarvis1openworldmultitaskagents}. Agents must reason about objects that may be out of view, plan sequences of actions that may take hundreds or thousands of steps to complete, and adapt to unexpected environmental changes. These challenges make Minecraft particularly valuable for testing the limits of current AI approaches and driving innovations in planning, perception, and learning algorithms.

Beyond technical challenges, Minecraft enables social interaction and emergent behaviors that are difficult to study in more constrained environments \cite{al2024projectsidmanyagentsimulations, yu2024minelandsimulatinglargescalemultiagent}. The MineLand simulator, for example, introduces features like limited multimodal senses and physical needs to create more ecologically valid social interactions among large groups of agents \cite{yu2024minelandsimulatinglargescalemultiagent}. This capability for modeling complex social dynamics makes Minecraft an ideal testbed for studying how artificial societies might form, evolve, and sustain themselves over time.

\subsection{Approaches to Agent Collaboration}

Agent collaboration approaches typically fall into two broad categories: task-oriented collaboration, which focuses on completing specific objectives, and social-oriented collaboration, which emphasizes emergent social structures and norms \cite{shu2024effectivegenaimultiagentcollaboration,qiu2024collaborativeintelligencepropagatingintentions}. Task-oriented approaches typically define explicit goals, roles, and protocols for interaction, prioritizing efficiency and measurable outcomes. Social-oriented approaches, in contrast, allow organizational structures and behavioral norms to emerge through repeated interactions, often leading to more adaptable but less predictable collaborative behaviors.

Coordination mechanisms in multi-agent systems include role assignment, task allocation, and conflict resolution strategies \cite{dong2024villageragentgraphbasedmultiagentframework,shu2024effectivegenaimultiagentcollaboration}. Role-based approaches assign specific responsibilities to agents based on their capabilities or position within an organizational structure. Task allocation methods distribute work among agents to optimize resource utilization and minimize conflicts. These mechanisms may be static (predetermined) or dynamic (adapting to changing conditions), with varying degrees of centralization in decision-making authority.

Communication strategies play a crucial role in effective collaboration, ranging from direct communication through message passing to indirect communication through environmental modifications \cite{shu2024effectivegenaimultiagentcollaboration,qiu2024collaborativeintelligencepropagatingintentions}. Direct communication allows for precise information sharing but requires protocols for when, what, and how to communicate. Indirect communication (stigmergy) relies on agents observing and responding to changes in the shared environment, often leading to emergent coordination patterns. The choice of communication strategy significantly impacts a system's scalability, robustness, and adaptability to novel situations.

Evaluation metrics for successful collaboration include task completion rates, resource efficiency, adaptability to changing conditions, and robustness to agent failures \cite{shu2024effectivegenaimultiagentcollaboration,qiu2024collaborativeintelligencepropagatingintentions, chen2024sagentsselforganizingagentsopenended}. Task-oriented systems typically prioritize metrics related to goal achievement and efficiency, such as time-to-completion and resource consumption. Social-oriented systems may emphasize metrics related to sustainability and resilience, such as the stability of social structures and adaptability to unexpected challenges. These different evaluation priorities reflect fundamental differences in how collaboration is conceptualized across approaches.

\section{Methodology}

\subsection{Case Study Selection}

The selection of Project Sid and TeamCraft as case studies for this comparative analysis was motivated by several key factors. First, both represent cutting-edge research published in 2024, ensuring relevance to the current state of multi-agent systems in Minecraft \cite{al2024projectsidmanyagentsimulations,long2024teamcraftbenchmarkmultimodalmultiagent}. Second, they exemplify contrasting approaches to multi-agent collaboration: Project Sid focuses on emergent social structures and civilizational processes, while TeamCraft emphasizes structured task completion through explicit coordination frameworks. This contrast provides a natural experiment for examining different paradigms of agent collaboration.

The comparison criteria for this analysis were developed based on key dimensions identified in the multi-agent systems literature, including organizational structure, communication patterns, learning mechanisms, and sustainability \cite{shu2024effectivegenaimultiagentcollaboration,qiu2024collaborativeintelligencepropagatingintentions, chen2024sagentsselforganizingagentsopenended}. These dimensions capture essential aspects of collaborative systems while allowing for meaningful comparison across different architectural approaches. Each dimension is further broken down into specific metrics and indicators that can be assessed based on the published descriptions and results of the two systems.

It is important to acknowledge the limitations of the case study approach. First, the analysis is constrained by the information available in published papers and associated materials, which may not capture all aspects of the systems' design and performance. Second, the two systems were developed with different primary objectives---Project Sid for studying emergent civilizational processes and TeamCraft for benchmarking multi-modal collaboration---making direct comparison challenging in some areas. Finally, as with any case study research, findings may not generalize to all multi-agent systems or environments beyond Minecraft.

\subsection{Analytical Framework}

The comparative analysis employs a structured framework examining four key dimensions: (1) Organizational structure, (2) Communication and coordination, (3) Learning and adaptation, and (4) Sustainability and scalability. Each dimension is analyzed through specific metrics drawn from multi-agent systems literature and organizational theory, providing a systematic basis for comparison.

Organizational structure is analyzed through the lens of hierarchy (the levels of authority and decision-making in the system), role specialization (how agents develop specialized functions), and adaptability (how the organization responds to changing conditions) \cite{Lhaksmana2018RoleBasedMF,chen2024sagentsselforganizingagentsopenended}. Project Sid's emergent hierarchies are contrasted with TeamCraft's task-oriented organization to understand how different structural approaches affect collaborative outcomes. This analysis draws on organizational theory concepts such as centralization, formalization, and departmentalization to characterize the systems' organizational features.

Agent autonomy and communication patterns are evaluated by examining communication frequency (how often agents exchange information), content (what information is shared), and effectiveness (how communication impacts task performance) \cite{shu2024effectivegenaimultiagentcollaboration,qiu2024collaborativeintelligencepropagatingintentions}. The analysis considers both direct communication through message passing and indirect communication through environmental modifications. Particular attention is paid to how communication patterns support or hinder coordination in different contexts, such as novel tasks or unexpected environmental changes.

Task allocation and sustainability are measured through metrics including efficiency (resource usage relative to task completion), adaptability (performance on novel tasks or with novel team compositions), and long-term viability (stability of performance over extended periods) \cite{shu2024effectivegenaimultiagentcollaboration,chen2024sagentsselforganizingagentsopenended}. These metrics provide insight into how well each system balances short-term performance goals with long-term sustainability concerns, a critical consideration for enduring multi-agent collaborations.

\subsection{Use of Generative AI}

This research employed generative AI tools throughout the research process, including literature review, analysis, and manuscript preparation. Specifically, LLMs were used to assist in identifying relevant literature, synthesizing findings across sources, and generating initial drafts of paper sections. This approach allowed for efficient processing of the extensive literature on multi-agent systems while maintaining focus on the specific research question.

The primary AI tool used was a LLM with research capabilities, which provided assistance in summarizing academic papers, identifying connections between concepts, and generating coherent prose based on specified content requirements. The model was provided with search results from academic databases and asked to summarize key findings, compare approaches, and identify research gaps. Human oversight was maintained throughout this process, with all AI-generated content reviewed, revised, and expanded based on direct examination of the original sources.

Ethical considerations in the use of generative AI included ensuring proper attribution of ideas to original sources, verifying factual claims against primary literature, and maintaining transparency about the role of AI in the research process. All content generated by AI was treated as a first draft requiring human verification and refinement rather than as finished work. This approach aligns with emerging best practices for AI-assisted academic writing, which emphasize human oversight and accountability.

The use of generative AI introduces certain limitations and potential biases, including the risk of propagating inaccuracies present in the training data, generating plausible-sounding but incorrect information, and prioritizing certain types of sources over others. To mitigate these risks, a systematic verification process was implemented, comparing AI-generated content against original sources and seeking additional sources where information appeared incomplete or potentially biased. The complete record of interactions with the AI system is included in the appendix to ensure full transparency about the research process.

% Generative AI (Perplexity and GPT-4) was used for literature search, synthesis, and drafting. All AI-generated content was reviewed and edited by the author. The full conversation history is included in the appendix for transparency, following best practices for responsible AI use \cite{lund2023chatgpt}.

\section{Case Study 1: Project Sid}

\subsection{Overview}

Project Sid represents a pioneering effort to study large-scale multi-agent interactions that mirror the full spectrum of civilizational processes. This research initiative simulates societies of 10 to over 1000 AI agents within a Minecraft environment, focusing on how these agents develop complex social structures and behaviors through extended interaction \cite{al2024projectsidmanyagentsimulations}. The project introduces the PIANO (Parallel Information Aggregation via Neural Orchestration) architecture, which enables agents to interact with humans and other agents in real-time while maintaining coherence across multiple output streams.

The PIANO architecture forms the technical foundation of Project Sid, providing a framework for agent perception, reasoning, and communication \cite{al2024projectsidmanyagentsimulations}. This architecture allows agents to process multimodal inputs, maintain internal states representing beliefs and goals, and generate appropriate actions in response to dynamic situations. A key innovation of PIANO is its ability to manage parallel information streams, enabling agents to simultaneously monitor environmental changes, track social interactions, and update their behavioral models accordingly.

Project Sid evaluates agent performance using civilizational benchmarks inspired by human history, measuring progress in areas such as role specialization, rule formation, and cultural transmission \cite{al2024projectsidmanyagentsimulations}. These benchmarks provide a structured way to assess whether agent societies are developing the key features that characterize human civilizations, moving beyond simple task completion metrics to capture the complexity of social evolution. The benchmarks include both quantitative measures (e.g., number of specialized roles) and qualitative assessments (e.g., coherence of cultural beliefs).

The simulation environment in Project Sid is carefully designed to support emergent social behaviors while maintaining computational tractability \cite{al2024projectsidmanyagentsimulations}. Agents exist in a shared Minecraft world where they can observe each other's actions, modify the environment, and engage in direct communication. The environment includes resources that must be gathered and transformed, creating natural interdependencies that encourage collaboration. This combination of shared space, limited resources, and communication capabilities creates the conditions necessary for complex social dynamics to emerge.

\subsection{Self-Organizing Social Systems}

One of the most striking findings of Project Sid is the spontaneous development of specialized roles among agents, mirroring the occupational specialization seen in human societies \cite{al2024projectsidmanyagentsimulations}. Without explicit programming of role definitions, agents naturally differentiate their behaviors based on environmental conditions, personal capabilities, and social feedback. This specialization emerges through a combination of reinforcement learning (where successful behaviors are repeated) and social learning (where agents imitate effective strategies observed in others). The resulting division of labor improves collective efficiency while creating interdependencies that strengthen social bonds.

The Project Sid agents demonstrate a remarkable ability to establish, adhere to, and modify collective rules that govern their interactions \cite{al2024projectsidmanyagentsimulations}. These rules emerge through repeated social interactions and feedback, gradually coalescing into informal norms and eventually more formalized governance structures. Agents learn to follow rules that benefit the collective, enforce these rules through social sanctions, and adapt rules when environmental or social conditions change. This process of rule formation and evolution parallels the development of social contracts and legal systems in human societies.

Cultural and religious transmission mechanisms represent another sophisticated aspect of Project Sid's agent societies \cite{al2024projectsidmanyagentsimulations}. Agents develop shared beliefs, values, and practices that spread through the population via direct communication and behavioral imitation. These cultural elements serve multiple functions: they coordinate collective activities, preserve useful knowledge, and strengthen group identity. In some simulations, belief systems resembling religious frameworks emerge, complete with ritual practices and mythological narratives that help agents make sense of their environment and relationships.

The emergent social behaviors observed in Project Sid extend beyond role specialization, rule formation, and cultural transmission to include complex patterns of cooperation, competition, and conflict resolution \cite{al2024projectsidmanyagentsimulations}. Agents form alliances based on shared interests, compete for limited resources when necessary, and develop mechanisms for resolving disputes without destructive conflict. These behaviors are not explicitly programmed but emerge from the agents' learning processes and repeated interactions, demonstrating how complex social dynamics can arise from relatively simple underlying mechanisms.

\subsection{Key Findings and Limitations}

Project Sid's most significant achievement is demonstrating that AI agents can autonomously develop civilizational milestones reminiscent of human societal evolution \cite{al2024projectsidmanyagentsimulations}. The agents successfully establish specialized roles, create and enforce social norms, develop cultural transmission mechanisms, and engage in complex collaborative behaviors---all without explicit programming of these capabilities. This emergent complexity suggests that fundamental principles of social organization may be discoverable through multi-agent learning, potentially offering insights into human social evolution as well as providing blueprints for artificial social systems.

Despite these achievements, Project Sid faces several challenges and limitations \cite{al2024projectsidmanyagentsimulations}. Computational constraints limit the scale and duration of simulations, potentially obscuring longer-term evolutionary dynamics. The agents' behavioral repertoire, while impressive, still falls short of human social complexity in areas such as moral reasoning, aesthetic appreciation, and abstract problem-solving. Additionally, the current implementation relies heavily on language model capabilities, making it difficult to disentangle which aspects of social behavior emerge from the agents' learning and which are implicit in the underlying models.

The implications of Project Sid for AI civilizations extend beyond academic interest to potential applications in organizational intelligence and human-AI collaboration \cite{al2024projectsidmanyagentsimulations}. The observed patterns of self-organization could inform the design of multi-agent systems for complex tasks such as disaster response, urban planning, or managing large-scale industrial processes. By understanding how agent societies naturally organize, developers could create more adaptable and resilient systems that require less explicit coordination overhead and adapt more readily to changing conditions.

Several aspects of Project Sid remain unexplored, pointing to potential areas for future research \cite{al2024projectsidmanyagentsimulations}. These include the long-term stability of agent societies over extended time periods, the impact of heterogeneous agent capabilities on social organization, and the potential for cross-cultural exchange between distinct agent societies. Additionally, more work is needed to understand how the principles observed in Minecraft environments might transfer to other virtual contexts or physical robotics applications, where different constraints and affordances may shape social dynamics in unexpected ways.

\section{Case Study 2: TeamCraft}

\subsection{Overview}

TeamCraft represents a significant advancement in multi-modal multi-agent benchmarking within the Minecraft environment, designed specifically to evaluate how effectively autonomous agents can collaborate on complex tasks through multi-modal understanding \cite{long2024teamcraftbenchmarkmultimodalmultiagent}. The benchmark features 55,000 task variants specified through multi-modal prompts, combining visual information, textual descriptions, and contextual cues to define goals and constraints. This diversity of tasks enables comprehensive testing of agents' ability to understand and execute collaborative activities across varying contexts and complexity levels.

The multi-modal prompts in TeamCraft span a wide range of task types, from simple resource gathering to complex construction projects requiring precise coordination \cite{long2024teamcraftbenchmarkmultimodalmultiagent}. These prompts combine visual elements (such as images of target structures or environments), textual instructions (describing goals and constraints), and contextual information (such as available resources or time limitations). This multi-modal approach tests agents' ability to integrate information across different perceptual channels, a critical capability for effective collaboration in rich environments.

A key innovation of TeamCraft is its extensive collection of procedurally-generated expert demonstrations for imitation learning \cite{long2024teamcraftbenchmarkmultimodalmultiagent}. These demonstrations provide examples of successful task completion, showing agents how to coordinate effectively to achieve goals. By varying demonstration parameters such as team size, resource availability, and environmental conditions, TeamCraft creates a rich dataset for training and evaluating collaborative behaviors. This approach supports data-efficient learning while maintaining diversity in the solution space.

TeamCraft implements carefully designed evaluation protocols to assess model generalization capabilities across several dimensions \cite{long2024teamcraftbenchmarkmultimodalmultiagent}. These protocols test agents' ability to generalize to novel goals (tasks not seen during training), novel scenes (environments with different layouts or resource distributions), and novel team compositions (varying the number of agents or their individual capabilities). This multi-dimensional evaluation provides a comprehensive picture of collaborative performance and helps identify specific generalization challenges for current models.

\subsection{Task-Oriented Coordination Frameworks}

Agents in TeamCraft must effectively understand and execute multi-modal tasks, requiring sophisticated perception and reasoning capabilities \cite{long2024teamcraftbenchmarkmultimodalmultiagent}. Visual inputs must be parsed to identify relevant environmental features, objects, and spatial relationships. Textual instructions must be interpreted to extract goal specifications, constraints, and sequential dependencies. These different information sources must then be integrated into a coherent task representation that guides planning and action selection. TeamCraft's framework evaluates how well agents perform this multi-modal understanding and how it impacts collaborative effectiveness.

Coordination mechanisms for collaborative tasks in TeamCraft include explicit role assignment and task decomposition strategies \cite{long2024teamcraftbenchmarkmultimodalmultiagent}. Agents must determine which subtasks to pursue individually versus collectively, how to allocate responsibilities based on agent capabilities and positions, and how to sequence activities to respect dependencies. These coordination decisions directly impact efficiency, with successful coordination reducing redundant efforts and ensuring proper sequencing of interdependent actions. TeamCraft provides metrics for evaluating coordination quality based on factors such as resource utilization, time efficiency, and goal achievement.

TeamCraft agents learn collaborative behaviors primarily through imitation learning from procedurally-generated demonstrations, supplemented by adaptation mechanisms for novel situations \cite{long2024teamcraftbenchmarkmultimodalmultiagent}. This approach allows agents to rapidly acquire effective coordination strategies from examples while developing the flexibility to adjust these strategies when faced with new contexts. The learning process encompasses both task-specific knowledge (how to complete particular goals) and generalizable collaboration principles (how to coordinate effectively regardless of the specific task).

Despite significant progress, TeamCraft agents face substantial challenges in generalizing to novel goals, scenes, and team compositions \cite{long2024teamcraftbenchmarkmultimodalmultiagent}. Performance typically degrades when agents encounter tasks that differ significantly from training examples, environments with unfamiliar layouts or resource distributions, or teams with different numbers of agents. These generalization challenges highlight the difficulty of developing truly adaptable collaborative systems and underscore the need for more robust approaches to knowledge transfer and adaptive coordination.

\subsection{Key Findings and Limitations}

TeamCraft's performance evaluation reveals both promising capabilities and significant limitations in current multi-agent collaboration models \cite{long2024teamcraftbenchmarkmultimodalmultiagent}. On familiar tasks in familiar environments with consistent team compositions, agents demonstrate impressive coordination capabilities, efficiently allocating subtasks and managing interdependencies. This performance suggests that contemporary models can successfully learn effective collaboration strategies from demonstrations when the testing conditions closely match training conditions.

The most significant challenge identified in TeamCraft is generalization across different dimensions of variation \cite{long2024teamcraftbenchmarkmultimodalmultiagent}. Performance degrades most severely when agents encounter novel goals requiring unfamiliar action sequences, suggesting limitations in how agents abstract collaborative principles from specific examples. Generalization to novel scenes shows moderate performance drops, while adapting to different team sizes presents particular difficulties, especially when scaling to larger teams than seen during training. These findings highlight the gap between current systems and the robust generalization capabilities needed for real-world applications.

Potential improvements to TeamCraft include architectural enhancements to better support cross-modal reasoning, more sophisticated communication protocols, and advanced meta-learning approaches \cite{long2024teamcraftbenchmarkmultimodalmultiagent}. Cross-modal attention mechanisms could improve integration of visual and textual information, while explicit communication channels could enhance coordination by allowing agents to share observations and intentions. Meta-learning approaches that explicitly optimize for adaptation to novel conditions might address the generalization challenges by encouraging agents to learn flexible collaboration strategies rather than task-specific solutions.

TeamCraft makes significant contributions to multi-agent benchmarking by establishing a comprehensive evaluation framework specifically focused on collaborative capabilities \cite{long2024teamcraftbenchmarkmultimodalmultiagent}. Prior benchmarks often emphasized individual performance or limited aspects of collaboration, whereas TeamCraft provides a holistic assessment across multiple dimensions of collaborative behavior. This comprehensive approach helps identify specific strengths and weaknesses in current models, guiding future research toward addressing the most critical limitations in multi-agent coordination.

\section{Comparative Analysis}

\subsection{Organizational Structures}

Project Sid and TeamCraft represent fundamentally different approaches to organizational structure in multi-agent systems. Project Sid employs an emergent organizational model where hierarchies, roles, and relationships develop through agent interactions without predefined structures \cite{al2024projectsidmanyagentsimulations}. Agents gradually differentiate into specialized roles based on their experiences and the needs of the collective, forming flexible hierarchies that evolve with changing conditions. This approach parallels how human societies naturally organize, with leadership and role specialization emerging from repeated social interactions rather than external design.

In contrast, TeamCraft implements a more structured, task-oriented organizational approach where roles and responsibilities are explicitly defined by task requirements \cite{long2024teamcraftbenchmarkmultimodalmultiagent}. Coordination frameworks provide clear guidelines for how agents should divide labor and sequence activities, creating more predictable but less adaptable organizational patterns. This approach resembles formal organizations in human society, where roles and processes are designed to optimize efficiency for specific objectives rather than evolving organically through social dynamics.

The decision-making processes in these systems reflect their organizational philosophies. Project Sid exhibits a mix of centralized and decentralized decision-making that emerges based on task complexity and social dynamics \cite{al2024projectsidmanyagentsimulations}. Leaders naturally emerge for certain activities, while other decisions remain distributed among agents. TeamCraft, however, tends toward more systematized decision processes guided by task decomposition and role assignments derived from expert demonstrations \cite{long2024teamcraftbenchmarkmultimodalmultiagent}. These different approaches create tradeoffs between adaptability and efficiency, with Project Sid's flexible decision-making allowing better adaptation to unexpected situations while TeamCraft's structured processes optimize performance on well-defined tasks.

Each organizational approach demonstrates distinct strengths and weaknesses when adapting to changing conditions \cite{al2024projectsidmanyagentsimulations,long2024teamcraftbenchmarkmultimodalmultiagent,chen2024sagentsselforganizingagentsopenended}. Project Sid's emergent structures show remarkable adaptability to environmental changes and novel social configurations, allowing the system to reorganize in response to new challenges without external intervention. However, this adaptability comes at the cost of efficiency, as emergent organization requires time to stabilize and may not optimize for specific task performance. TeamCraft's structured approach enables rapid mobilization for familiar tasks but struggles with novel contexts that require organizational reconfiguration, showing diminished performance when faced with unexpected challenges or team composition changes.

\subsection{Communication and Coordination}

Communication mechanisms differ significantly between Project Sid and TeamCraft, reflecting their contrasting organizational philosophies. Project Sid allows communication patterns to emerge naturally through agent interactions, with both direct communication (explicit message passing) and indirect communication (environmental modifications) developing as agents learn effective interaction strategies \cite{al2024projectsidmanyagentsimulations}. This emergent communication evolves to meet the needs of the collective, becoming more structured and specialized as social complexity increases.

TeamCraft, in contrast, implements more formalized communication protocols derived from expert demonstrations, with clear patterns for information sharing and coordination signaling \cite{long2024teamcraftbenchmarkmultimodalmultiagent,qiu2024collaborativeintelligencepropagatingintentions}. These protocols are optimized for task efficiency, ensuring that agents share precisely the information needed for successful collaboration without unnecessary communication overhead. While effective for familiar tasks, this approach may limit adaptability when novel situations require communication patterns not present in the training demonstrations.

Task allocation strategies similarly reflect the systems' divergent approaches to organization. Project Sid develops allocation mechanisms through trial and error as agents learn which distributions of responsibilities lead to collective success \cite{al2024projectsidmanyagentsimulations}. This emergent allocation adapts dynamically to changing circumstances and agent capabilities, allowing the system to reorganize in response to new challenges. TeamCraft implements more explicit allocation strategies based on task structure analysis, optimizing assignments to minimize completion time and resource usage \cite{long2024teamcraftbenchmarkmultimodalmultiagent,dong2024villageragentgraphbasedmultiagentframework}. This approach yields high efficiency for well-understood tasks but may struggle with novel situations requiring creative reallocation of responsibilities.

The different communication and coordination approaches create fundamental trade-offs in system performance \cite{al2024projectsidmanyagentsimulations,long2024teamcraftbenchmarkmultimodalmultiagent,qiu2024collaborativeintelligencepropagatingintentions,dong2024villageragentgraphbasedmultiagentframework}. Project Sid's emergent strategies develop greater robustness to unexpected changes and novel situations, as the system can adapt its communication and coordination patterns through continued learning. However, this adaptability comes at the cost of initial efficiency, as effective patterns must evolve through experience rather than being optimized from the start. TeamCraft's structured approaches deliver superior performance on familiar tasks through optimized protocols but show limited flexibility when facing novel challenges that require communication or coordination strategies not present in training data.

\subsection{Learning and Adaptation}

The learning mechanisms employed by Project Sid and TeamCraft represent distinct approaches to knowledge acquisition and behavioral adaptation. Project Sid primarily utilizes reinforcement learning combined with social learning, allowing agents to discover effective behaviors through environmental feedback while also learning from observing other agents \cite{al2024projectsidmanyagentsimulations}. This approach enables open-ended learning where new capabilities can emerge without being explicitly encoded in training data, facilitating the development of novel social structures and interaction patterns.

TeamCraft relies more heavily on imitation learning from expert demonstrations, supplemented with fine-tuning through reinforcement learning \cite{long2024teamcraftbenchmarkmultimodalmultiagent,zhao2024steveseriesstepbystepconstruction}. This approach allows agents to rapidly acquire effective coordination strategies by observing successful examples, learning both task-specific knowledge and generalizable collaboration principles. The structured learning process produces more predictable performance on tasks similar to the demonstrations but may limit discovery of novel strategies that weren't present in the training data.

Adaptation to novel situations reveals significant differences between the two systems. Project Sid demonstrates stronger adaptation to completely novel scenarios due to its emphasis on foundational learning mechanisms rather than task-specific knowledge \cite{al2024projectsidmanyagentsimulations}. Agents can apply basic principles of social organization and resource management to unfamiliar challenges, gradually developing effective responses through continued interaction and feedback. TeamCraft shows better immediate performance on variations of familiar tasks but struggles with fundamentally novel challenges that require significant departures from demonstrated behaviors \cite{long2024teamcraftbenchmarkmultimodalmultiagent}.

Both systems face challenges in knowledge transfer across domains and tasks, though for different reasons \cite{al2024projectsidmanyagentsimulations,long2024teamcraftbenchmarkmultimodalmultiagent,zhao2024steveseriesstepbystepconstruction}. Project Sid's emergent knowledge is often implicit in agent behaviors rather than explicitly represented, making it difficult to transfer learning directly from one context to another. TeamCraft's more structured representations enable clearer knowledge transfer within similar task families but may create overfitting to particular task structures that hinders transfer to significantly different domains. These different learning approaches highlight the tension between flexibility and efficiency in knowledge acquisition for multi-agent systems.

\subsection{Sustainability and Scalability}

The long-term viability of agent collaborations differs markedly between Project Sid and TeamCraft. Project Sid demonstrates superior sustainability through its self-organizing social structures, which continuously adapt to changing conditions and can recover from perturbations without external intervention \cite{al2024projectsidmanyagentsimulations,yu2024minelandsimulatinglargescalemultiagent}. The development of cultural transmission mechanisms allows knowledge to persist across time, while emergent governance structures help maintain social cohesion even as individual agents come and go. These features create resilient social systems capable of enduring over extended periods despite changing environmental conditions.

TeamCraft shows stronger performance in short-term task execution but may face challenges in long-term sustainability due to its reliance on predefined coordination frameworks \cite{long2024teamcraftbenchmarkmultimodalmultiagent}. Without continuous adaptation mechanisms, these frameworks may become less effective as conditions drift from their original design parameters. The system lacks the self-organizing capabilities that would allow it to reorganize in response to significant environmental changes or agent population shifts, potentially leading to performance degradation over extended periods.

Scaling to larger agent populations presents distinct challenges for each system. Project Sid demonstrates impressive scalability, successfully operating with populations ranging from 10 to over 1000 agents \cite{al2024projectsidmanyagentsimulations}. The emergent organizational structures naturally accommodate growing populations through increasing specialization and hierarchical organization, mirroring how human societies scale. This natural scalability comes at the cost of increasing organizational complexity and potential inefficiencies as coordination requirements grow more complex.

TeamCraft faces greater challenges in scaling beyond the team sizes present in its training data \cite{long2024teamcraftbenchmarkmultimodalmultiagent,yu2024minelandsimulatinglargescalemultiagent}. Performance tends to degrade when the number of agents increases significantly beyond what was demonstrated, suggesting limitations in how the coordination frameworks generalize to larger teams. This scaling challenge highlights a broader issue with task-oriented approaches: their optimization for specific contexts may create implicit assumptions about team size and composition that limit generalization to significantly different configurations.

Resource efficiency presents different tradeoffs in the two systems. Project Sid initially shows lower efficiency as social structures are forming but develops increasingly efficient resource utilization as specialization and coordination mechanisms mature \cite{al2024projectsidmanyagentsimulations}. TeamCraft demonstrates higher immediate efficiency on familiar tasks but may use resources suboptimally when facing novel challenges that require adaptation \cite{long2024teamcraftbenchmarkmultimodalmultiagent}. These patterns highlight the fundamental tension between optimization for known conditions versus adaptability to changing circumstances---a key consideration for sustainable multi-agent systems.

\section{Discussion}

\subsection{Theoretical Implications}

The comparison of Project Sid and TeamCraft yields important insights for multi-agent systems theory, particularly regarding the relationship between organizational structure and adaptive capacity. The emergent, self-organizing approach of Project Sid demonstrates how complex coordination can arise without explicit design, suggesting that fundamental principles of social organization may be discoverable through agent learning rather than requiring human engineering \cite{al2024projectsidmanyagentsimulations,chen2024sagentsselforganizingagentsopenended}. This finding supports theories of emergent complexity in multi-agent systems and suggests that allowing structural flexibility may be crucial for developing truly adaptive collaborative systems.

The contrast between these approaches contributes significantly to our understanding of artificial societies, revealing parallels with human social systems while also highlighting uniquely computational aspects of agent societies \cite{al2024projectsidmanyagentsimulations,long2024teamcraftbenchmarkmultimodalmultiagent}. Like human societies, agent groups in Project Sid develop specialized roles, governance structures, and cultural transmission mechanisms that enhance collective capabilities. Unlike human societies, however, these developments occur on dramatically accelerated timescales and may follow different evolutionary trajectories due to the agents' computational nature. These similarities and differences provide a novel lens for studying social organization more broadly.

The organizational patterns observed in these systems connect directly to human organizational behavior theories, offering insights for designing effective human-AI collaborative systems \cite{chen2024sagentsselforganizingagentsopenended,Lhaksmana2018RoleBasedMF}. Project Sid's emergent leadership and role specialization mirror organic organizational development in human groups, while TeamCraft's structured coordination frameworks resemble formal organizational designs. These parallels suggest that principles from organizational theory---such as the importance of matching structure to environmental uncertainty---may apply equally to artificial agent societies, informing how we design systems for different collaborative contexts.

The development of self-organizing AI societies raises profound philosophical questions about autonomy, purpose, and value alignment in artificial systems \cite{al2024projectsidmanyagentsimulations}. As agent societies develop their own norms, governance structures, and cultural practices, they begin to exhibit a form of collective autonomy that transcends individual agent programming. This emergence challenges traditional notions of AI system design, where behaviors are explicitly encoded rather than emergently developed. Understanding how values and purposes evolve in these emergent systems becomes crucial for ensuring alignment with human priorities, particularly as agent societies grow in complexity and independence.

\subsection{Practical Applications}

The findings from comparing Project Sid and TeamCraft have significant implications for real-world applications of multi-agent systems. Organizational design could benefit from hybrid approaches that combine TeamCraft's efficiency in well-structured domains with Project Sid's adaptability in dynamic environments \cite{long2024teamcraftbenchmarkmultimodalmultiagent,al2024projectsidmanyagentsimulations}. For example, disaster response systems might employ predefined coordination protocols for routine operations while maintaining the capacity for emergent reorganization when facing unprecedented situations. This balanced approach could optimize performance across varying conditions while maintaining robustness to unexpected challenges.

Insights from Minecraft-based systems could transfer to other domains including physical robotics, business process management, and social network governance \cite{dong2024villageragentgraphbasedmultiagentframework, zhao2024hierarchicalautoorganizingopenendedmultiagent}. The coordination mechanisms observed in TeamCraft could inform task allocation algorithms for warehouse robotics or manufacturing systems, while Project Sid's emergent governance structures might inspire approaches to community management in online platforms. The abstract principles of effective collaboration---such as balancing specialization with flexibility---transcend the specific Minecraft environment and can guide system design across diverse applications.

Several design principles for effective multi-agent systems emerge from this comparison. First, incorporating both structured coordination frameworks and emergent organizational capabilities creates systems that can optimize for efficiency while maintaining adaptability \cite{chen2024sagentsselforganizingagentsopenended,qiu2024collaborativeintelligencepropagatingintentions}. Second, developing robust communication protocols that balance information sharing with communication overhead enhances coordination across varying contexts \cite{shu2024effectivegenaimultiagentcollaboration,qiu2024collaborativeintelligencepropagatingintentions}. Third, supporting role flexibility while encouraging specialization allows systems to balance expertise development with adaptive capacity \cite{al2024projectsidmanyagentsimulations,long2024teamcraftbenchmarkmultimodalmultiagent}. Fourth, implementing cultural transmission mechanisms enables knowledge preservation across time while facilitating system-wide adaptation to changing conditions \cite{al2024projectsidmanyagentsimulations}.

Deploying multi-agent systems in real-world contexts raises important ethical considerations regarding transparency, accountability, and human oversight \cite{chen2024sagentsselforganizingagentsopenended}. As these systems develop increasingly autonomous organizational capabilities, ensuring they remain aligned with human values and priorities becomes crucial. Mechanisms for explaining emergent behaviors, maintaining human control over critical decisions, and verifying system performance across diverse conditions must be incorporated into system design from the outset. These ethical safeguards are particularly important as multi-agent systems scale to include larger numbers of agents with more complex interaction patterns.

\section{Conclusion and Future Work}

This comparative analysis of Project Sid and TeamCraft reveals fundamental differences in how multi-agent systems approach collaboration in Minecraft environments. Project Sid's emergent, socially-oriented approach produces adaptable, self-organizing agent societies capable of developing specialized roles, governance structures, and cultural transmission mechanisms \cite{al2024projectsidmanyagentsimulations}. In contrast, TeamCraft's structured, task-oriented approach enables efficient coordination through explicit frameworks optimized for specific collaborative challenges \cite{long2024teamcraftbenchmarkmultimodalmultiagent}. These different approaches create distinct strengths and limitations, with Project Sid excelling in adaptability and long-term sustainability while TeamCraft demonstrates superior efficiency on well-defined tasks.

Addressing our research question directly, Project Sid's self-organizing social systems achieve sustainable multi-agent collaboration through emergent structures that continuously adapt to changing conditions, enabling long-term viability despite environmental fluctuations and agent turnover \cite{al2024projectsidmanyagentsimulations}. TeamCraft's task-oriented coordination frameworks achieve sustainability through efficient resource utilization and precise task execution but may struggle with adaptation to significant environmental changes or team composition shifts \cite{long2024teamcraftbenchmarkmultimodalmultiagent}. The optimal approach depends on the specific collaborative context: dynamic, unpredictable environments favor emergent organization, while stable, well-defined domains benefit from structured coordination.

This study has several limitations that should be acknowledged. First, the analysis relies primarily on published descriptions rather than direct experimentation with the systems, potentially missing nuances in implementation and performance. Second, the two systems were developed with different primary objectives---Project Sid for studying emergent civilizational processes and TeamCraft for benchmarking multi-modal collaboration---making direct comparison challenging in some areas. Third, both systems are relatively recent developments, limiting our understanding of their long-term evolution and performance.

Future research should explore hybrid approaches that combine the strengths of both paradigms, potentially developing systems that employ structured coordination for efficiency while maintaining capacity for emergent reorganization when facing novel challenges \cite{qiu2024collaborativeintelligencepropagatingintentions,chen2024sagentsselforganizingagentsopenended}. Longitudinal studies examining how agent societies evolve over extended periods could provide deeper insights into sustainability and adaptive capacity. Integration with human societies represents another promising direction, investigating how artificial agent groups might collaborate with human teams while maintaining alignment with human values and priorities. Finally, exploring how insights from Minecraft-based systems transfer to physical robotics and real-world organizational contexts remains an important frontier for multi-agent systems research.

%%%%%%%%%%%%%%%%%%%%%%%%%%%%%%%%%%%%%%%%%%%%%%%%%%%%%%%%%%%%%
%BIBLIOGRAPHY
%%%%%%%%%%%%%%%%%%%%%%%%%%%%%%%%%%%%%%%%%%%%%%%%%%%%%%%%%%%%%

\clearpage
\bibliographystyle{apacite}
% \setlength{\bibleftmargin}{-1cm}
% \setlength{\bibindent}{-1cm}
\bibliography{literature.bib}


%%%%%%%%%%%%%%%%%%%%%%%%%%%%%%%%%%%%%%%%%%%%%%%%%%%%%%%%%%%%%
%APPENDICES
%%%%%%%%%%%%%%%%%%%%%%%%%%%%%%%%%%%%%%%%%%%%%%%%%%%%%%%%%%%%%

\clearpage
\appendix
%\renewcommand*{\thesection}{\Alph{section}}\textbf{}

% APPENDIX A
\section{Appendix}

\subsection{AI Conversation History}

\lstinputlisting[language={}, label={lst:prompt}]{includes/perplexity.md}


\end{document}